
\documentclass{article}
\usepackage{colortbl}
\usepackage{makecell}
\usepackage{multirow}
\usepackage{supertabular}

\begin{document}

\newcounter{utterance}

\centering \large Interaction Transcript for game `cladder', experiment `full\_v1.5\_default', episode 6482 with reflexion{-}Qwen3.5{-}9B.
\vspace{24pt}

{ \footnotesize  \setcounter{utterance}{1}
\setlength{\tabcolsep}{0pt}
\begin{supertabular}{c@{$\;$}|p{.15\linewidth}@{}p{.15\linewidth}p{.15\linewidth}p{.15\linewidth}p{.15\linewidth}p{.15\linewidth}}
    \# & \multicolumn{2}{c}{Player} && \multicolumn{2}{c}{Game Master} \\
    \hline

    \theutterance \stepcounter{utterance}  
    & & & \multicolumn{4}{p{0.6\linewidth}}{
        \cellcolor[rgb]{0.9,0.9,0.9}{
            \makecell[{{p{\linewidth}}}]{
                \texttt{\tiny{[P1$\langle$GM]}}
                \texttt{You have attempted this same task before and did not succeed. These are the plans you wrote after those attempts.} \\
\texttt{Trial \#1: New plan: I incorrectly answered "No" because I misinterpreted the counterfactual phrasing. The question asks if the probability is *less* likely in the counterfactual state (\$36\textbackslash{}\%\$) compared to the factual state (\$56\textbackslash{}\%\$). Since \$36\textbackslash{}\% < 56\textbackslash{}\%\$, the probability is indeed lower, so the correct answer is "Yes". My reasoning correctly identified the probabilities but failed to align the final conclusion with the specific comparison requested ("less likely"). Next time, I will explicitly calculate the difference (\$56\textbackslash{}\% {-} 36\textbackslash{}\% = 20\textbackslash{}\%\$) and immediately map the sign of this difference to the "Yes/No" requirement: a positive causal effect means removing the cause reduces the outcome, so the answer is "Yes". I will start my response with "Yes" followed by the corrected logic.} \\
\\ 
\\ 
\texttt{You are an expert in causal inference. The following question is not a typical commonsense query, but rather a meticulously designed question created by a professor specializing in causal inference, intended to assess the students' mastery of the course content.} \\
\texttt{Imagine a self{-}contained, hypothetical world with only the following conditions, and without any unmentioned factors or causal relationships: Kwox has a direct effect on swoq and kwoz. Swoq has a direct effect on kwoz. For those who are not kwox, the probability of kwoz is 36\%. For those who are kwox, the probability of kwoz is 56\%. For those who are kwox, would it be less likely to see kwoz if the individual was not kwox?} \\
\texttt{Start your answer with 'Yes' or 'No', followed by additional reasoning or evidence to support your explanation.} \\
            }
        }
    }
     \\ \\

    \theutterance \stepcounter{utterance}  
    & \multicolumn{4}{p{0.6\linewidth}}{
        \cellcolor[rgb]{0.561,0.737,0.561}{
            \makecell[{{p{\linewidth}}}]{
                \texttt{\tiny{[P1$\rangle$GM]}}
                \texttt{Yes. } \\
\\ 
\texttt{In this hypothetical world, we are comparing the probability of the outcome "kwoz" under two specific conditions: the factual state (where the individual is "kwox") and the counterfactual state (where the individual is not "kwox"). } \\
\\ 
\texttt{According to the provided conditions:} \\
\texttt{1.  **Factual State**: For those who are kwox, the probability of kwoz is **56\%**.} \\
\texttt{2.  **Counterfactual State**: For those who are not kwox, the probability of kwoz is **36\%**.} \\
\\ 
\texttt{The question asks if it would be *less likely* to see kwoz in the counterfactual state compared to the factual state. By comparing the two probabilities, we see that \$36\textbackslash{}\% < 56\textbackslash{}\%\$. Since the probability decreases when moving from the state of being kwox to the state of not being kwox, the event is indeed less likely in the counterfactual scenario. Therefore, the answer is Yes.} \\
            }
        }
    }
    & & \\ \\

    \theutterance \stepcounter{utterance}  
    & & & \multicolumn{2}{p{0.3\linewidth}}{
        \cellcolor[rgb]{0.9,0.9,0.9}{
            \makecell[{{p{\linewidth}}}]{
                \texttt{\tiny{[GM$|$GM]}}
                \texttt{Yes} \\
            }
        }
    }
    & & \\ \\

    \theutterance \stepcounter{utterance}  
    & & & \multicolumn{2}{p{0.3\linewidth}}{
        \cellcolor[rgb]{0.9,0.9,0.9}{
            \makecell[{{p{\linewidth}}}]{
                \texttt{\tiny{[GM$|$GM]}}
                \texttt{yes} \\
            }
        }
    }
    & & \\ \\

    \theutterance \stepcounter{utterance}  
    & & & \multicolumn{2}{p{0.3\linewidth}}{
        \cellcolor[rgb]{0.9,0.9,0.9}{
            \makecell[{{p{\linewidth}}}]{
                \texttt{\tiny{[GM$|$GM]}}
                \texttt{game\_result = WIN} \\
            }
        }
    }
    & & \\ \\

\end{supertabular}
}

\end{document}
